% Module 1 section file. Included by ../module1_stellarator_equilibrium_geometry.tex.
\section{Magnetic flux coordinates}
\label{sec:magnetic_coordinates}

\subsection{Coordinate mapping}
Let
\begin{equation}
(u^1,u^2,u^3)=(s,\theta,\zeta)
\end{equation}
be curvilinear coordinates. The physical position is a mapping
\begin{equation}
\mathbf x=\mathbf x(s,\theta,\zeta).
\label{eq:coordinate_mapping}
\end{equation}
The covariant basis vectors are
\begin{equation}
\mathbf e_i=\frac{\partial\mathbf x}{\partial u^i},
\end{equation}
and the reciprocal, or contravariant, basis vectors are
\begin{equation}
\mathbf e^i=\nabla u^i.
\end{equation}
They satisfy
\begin{equation}
\mathbf e_i\cdot\mathbf e^j=\delta_i^{\ j},
\end{equation}
where $\delta_i^{\ j}$ is the Kronecker delta, equal to one when $i=j$ and zero otherwise.

\subsection{Metric tensors}
The covariant metric tensor is
\begin{equation}
g_{ij}=\mathbf e_i\cdot\mathbf e_j,
\label{eq:covariant_metric}
\end{equation}
and the contravariant metric tensor is
\begin{equation}
g^{ij}=\mathbf e^i\cdot\mathbf e^j.
\label{eq:contravariant_metric}
\end{equation}
The matrices $[g_{ij}]$ and $[g^{ij}]$ are inverses. They convert coordinate derivatives and components into physical lengths and dot products.

For example, the squared differential distance is
\begin{equation}
d\ell^2=g_{ij}\,du^i du^j,
\label{eq:line_element}
\end{equation}
where repeated indices are summed from 1 to 3.

\subsection{Jacobian and volume element}
The coordinate Jacobian is
\begin{equation}
\sqrt{g}=\mathbf e_s\cdot(\mathbf e_\theta\times\mathbf e_\zeta).
\label{eq:jacobian}
\end{equation}
The physical volume element is
\begin{equation}
dV=\sqrt{g}\,ds\,d\theta\,d\zeta.
\label{eq:volume_element}
\end{equation}
The sign of $\sqrt{g}$ depends on coordinate orientation. A production interface should either enforce a positive orientation or store the sign convention explicitly.

\subsection{Contravariant representation of a tangent, divergence-free field}
A particularly useful straight-field-line representation is
\begin{equation}
\boxed{
\mathbf B
=\nabla\psi\times\nabla\theta
+\iota(\psi)\,\nabla\zeta\times\nabla\psi.
}
\label{eq:B_contravariant}
\end{equation}
Here $\psi$ is a toroidal-flux-like surface coordinate. Equation~\eqref{eq:B_contravariant} automatically gives
\begin{equation}
\mathbf B\cdot\nabla\psi=0
\end{equation}
and, provided $\iota$ depends only on $\psi$, satisfies $\nabla\cdot\mathbf B=0$.

Using the Jacobian identity
\begin{equation}
(\nabla\psi\times\nabla\theta)\cdot\nabla\zeta=\frac{1}{\sqrt{g}},
\end{equation}
we find
\begin{align}
\mathbf B\cdot\nabla\zeta&=\frac{1}{\sqrt{g}},\\
\mathbf B\cdot\nabla\theta&=\frac{\iota}{\sqrt{g}}.
\end{align}
Therefore
\begin{equation}
\frac{d\theta}{d\zeta}
=\frac{\mathbf B\cdot\nabla\theta}{\mathbf B\cdot\nabla\zeta}
=\iota,
\end{equation}
which explains the phrase \textbf{straight-field-line coordinates}: in the $(\theta,\zeta)$ plane, field lines are straight lines of slope $\iota$.

\subsection{Covariant representation and Boozer coordinates}
Any vector can also be expanded in gradients:
\begin{equation}
\mathbf B=B_s\nabla s+B_\theta\nabla\theta+B_\zeta\nabla\zeta.
\end{equation}
In Boozer coordinates the covariant representation can be written
\begin{equation}
\mathbf B
=K(s,\theta,\zeta)\nabla s
+I(s)\nabla\theta
+G(s)\nabla\zeta.
\label{eq:B_boozer_covariant}
\end{equation}
The functions $I(s)$ and $G(s)$ are flux functions related to net toroidal and poloidal currents, while $K$ may vary over a surface. Exact multiplicative and sign conventions depend on how flux and angles are normalized. This is why a geometry file must store metadata rather than only arrays.

Boozer coordinates are useful because the magnetic-field-line trajectories are straight and the covariant angular components are flux functions. Guiding-center orbit equations and quasisymmetry are especially transparent in this representation. Detailed treatments of flux coordinates and the original Boozer construction are given in Refs.~\cite{dHaeseleer1991,boozer1981}.

\subsection{Hamada coordinates}
Hamada coordinates are another straight-field-line system. They are constructed so that both magnetic field lines and current lines have simple contravariant representations and the Jacobian is a flux function under suitable conditions. Boozer and Hamada coordinates serve different analytical and numerical purposes; neither is simply a geometric angle system.

\subsection{Coordinate singularity at the magnetic axis}
At $s=0$, all poloidal angles refer to the same physical magnetic axis. The coordinate system is therefore singular in the same way that polar angle is undefined at the origin. This is a coordinate singularity, not necessarily a physical singularity. Numerical representations must impose regularity conditions so that physical fields remain finite and single-valued near the axis.

